\section{Results}

\subsection{Model fitness is separable from reward}

The first analysis separates model fitness from realized reward. If the beta-derived precision signal were functioning as cached value, it would be expected to covary with payoff at least as strongly as with predictive reliability. A model-fitness account instead expects precision to remain high for partners whose behavior is predictable, even when interaction with those partners is not especially rewarding.

Across 30 seeds per variant and 200 rounds per seed, affective precision did not improve payoff: total payoff was $534.6$ for the precision-modulated agent and $542.1$ for the no-affect baseline, with a mean difference of $-7.53$ and bootstrap CI $[-27.88,13.82]$. The absence of a payoff advantage is useful because it makes the reliability analysis a cleaner test of what the signal tracks.

The beta-derived precision signal tracked predictive reliability more strongly than reward. Across partner-seed units, $|\mathrm{corr}(\precision,\mathrm{surprise})|=0.701$, while $|\mathrm{corr}(\precision,\mathrm{payoff})|=0.419$. The corresponding seed-level reliability-over-reward summary was $+0.096$ with bootstrap CI $[0.027,0.164]$. This pattern is consistent with the model-fitness interpretation: the affective signal estimates confidence in the partner model, not partner value.

\begin{figure}[H]
\centering
\includegraphics[width=\linewidth]{fig_model_fitness_beta_reward_divergence.png}
\caption{Model-fitness measurement. In the 30-seed reliability-versus-reward analysis, partner-local precision tracks predictive surprise more strongly than realized payoff, while precision modulation has no total-payoff advantage.}
\label{fig:model-fitness}
\end{figure}

\subsection{Precision gates deployment in open policy regimes}

Partner-local precision can only change behavior when several policies remain live enough for precision to matter. In a saturated regime, the best policy may already dominate the posterior; changing softmax precision then changes little. In the open graded regime, using five seeds per variant and 200 rounds per seed, partner-local precision was associated with higher payoff and lower policy entropy. The precision-modulated agent reached total payoff $1884.6$ with mean policy entropy $7.94$, while the no-affect or tracked-only baseline reached $1859.4$ with entropy $8.80$. The difference was therefore $+25.2$ payoff and $-0.86$ entropy for precision modulation.

The tracked-only lesion separates belief estimation from action selection. In the same open regime, the precision-modulated agent had joint belief accuracy $0.336$ and stance accuracy $0.816$, while the tracked-only lesion had joint accuracy $0.339$ and stance accuracy $0.855$. Belief measurements were therefore similar or slightly favored the lesion, but policy entropy and payoff favored the precision pathway. The difference appears at deployment: the signal changes how beliefs are acted upon rather than simply improving belief accuracy.

\begin{figure}[H]
\centering
\includegraphics[width=\linewidth]{fig_deployment_social_summary.png}
\caption{Deployment and social choice measurements. In the open graded regime, local beta coupling changes payoff and policy entropy even when belief measurements remain similar. In the partner-choice regime, selection rates move before total payoff separates.}
\label{fig:deployment-social}
\end{figure}

\subsection{Partner choice changes before payoff}

The partner-choice regime shows how partner-local precision shapes approach, avoidance, and probing. Payoff is a blunt measurement here because partner selection can change before total payoff separates. Across five seeds and 200 rounds per seed, payoff was nearly flat: $393.6$ for precision modulation versus $393.2$ for no affect. The policy distribution nevertheless changed. Mean policy entropy was $3.99$ with precision modulation and $4.83$ without it. Selection shifted toward the first partner ($0.348$ to $0.366$) and away from the third partner ($0.162$ to $0.138$). Within this bounded evidence, partner-local precision changed approach and avoidance before total payoff became informative.

\subsection{Abrupt betrayal exposes a boundary condition}

Abrupt betrayal reveals a boundary condition in which the same precision channel can be harmful. A partner's stance changes suddenly, creating a mismatch between the agent's accumulated model and the partner's new behavior. Across 30 seeds per variant and 120 rounds per seed, the precision-modulated agent obtained lower whole-run payoff than no-affect or tracked-only baselines: $1136.1$ versus $1172.1$, with mean difference $-36.08$, bootstrap CI $[-63.65,-10.93]$, and pairwise $p=0.0169$. Policy entropy was lower under precision modulation ($8.38$ versus $8.74$), showing that the channel was behaviorally active.

Post-switch measurements explain why this is not affective recovery. Precision modulation reduced mean reencounters with the switched partner ($4.4$ versus $6.1$) and lowered reencounter selection rate ($0.049$ versus $0.067$), but it did not improve the safety of returning: payoff on reencounter was $8.76$ versus $8.91$, and wrong-type rate on reencounter was $0.237$ versus $0.165$. The agent returned less often, but when it returned, it was not safer or better calibrated. Abrupt betrayal therefore produced precision-driven misdeployment: decisions became sharper before the partner model had recalibrated enough.

\begin{figure}[H]
\centering
\includegraphics[width=\linewidth]{fig_betrayal_boundary_summary.png}
\caption{Betrayal boundary condition. Abrupt betrayal makes partner-local precision behaviorally visible, but not safer: precision modulation lowers entropy and changes reallocation while reducing whole-run payoff.}
\label{fig:betrayal-boundary}
\end{figure}

\subsection{Shock shape matters more than generic caution}

The precision-sensitivity analysis tests whether betrayal failure is simply excessive precision gain. In the abrupt sensitivity run, the no-affect/tracked-only baseline had payoff $1153.6$ and entropy $8.80$. Default precision modulation reached payoff $1140.4$ and entropy $8.35$ ($p=0.370$ against no affect). A combined-caution variant raised entropy to $9.27$ but lowered payoff to $1115.0$ ($p=0.003$ against no affect).

The gradual betrayal run changed the outcome. The no-affect/tracked-only baseline had payoff $1148.9$ and entropy $8.76$, while default precision modulation nearly matched it with payoff $1147.6$, entropy $8.38$, and $p=0.906$. Combined caution again reduced payoff ($1118.3$, $p=0.005$). Generic caution therefore does not rescue abrupt betrayal. The more informative boundary is shock shape: default partner-local precision is most harmful when a partner changes abruptly.

\begin{figure}[H]
\centering
\includegraphics[width=\linewidth]{fig_shock_shape_summary.png}
\caption{Shock-shape sensitivity. Default local beta has a larger payoff penalty under abrupt betrayal than gradual betrayal, while a generic caution variant raises entropy but remains costly in both regimes.}
\label{fig:shock-shape}
\end{figure}

\subsection{Perturbations separate precision dynamics before interpretation}

The perturbation runs are not clinical validation; they are a computational phenotyping check on the beta pathway. In open graded interaction, the low-gain variant compressed beta movement (mean beta range $0.149$ versus $0.950$ for default local beta), while high-gain and cautious-prior variants expanded it ($1.318$ and $1.435$). In betrayal, the same qualitative ordering remained visible: low gain stayed blunted ($0.132$), while high gain and cautious prior remained more volatile ($1.114$ and $1.091$). These differences are useful because they appear directly in the precision dynamics rather than only in payoff.

Behavioral measurements were less one-dimensional. Action churn stayed high across variants, and payoff did not rank variants by a simple ``more beta movement is worse'' rule. That pattern is the point of treating these as phenotypes rather than diagnoses: the perturbations define different beta dynamics, but the behavioral consequence depends on the regime in which those dynamics are deployed.

\begin{figure}[H]
\centering
\includegraphics[width=\linewidth]{fig_phenotype_dynamics_summary.png}
\caption{Precision-dynamics perturbation phenotypes. Low-gain, high-gain, and cautious-prior variants separate clearly in beta movement, while action churn and payoff depend on the task regime. These variants are computational perturbations, not clinical models.}
\label{fig:phenotype-dynamics}
\end{figure}
